\documentclass[10pt,a4paper]{article}
\usepackage[margin=1.7cm,top=1.5cm,bottom=1.5cm]{geometry}
\usepackage{amsmath,amssymb,booktabs,xcolor,titlesec,enumitem,array}
\usepackage[T1]{fontenc}
\definecolor{ac}{HTML}{1F5FBF}
\definecolor{bx}{HTML}{F3F6FB}
\titleformat{\section}{\large\bfseries\color{ac}}{\thesection}{0.6em}{}[\vspace{-0.55em}\color{ac}\rule{\linewidth}{0.5pt}]
\titlespacing*{\section}{0pt}{9pt}{5pt}
\setlength{\parindent}{0pt}
\setlength{\parskip}{3pt}
\setlist[itemize]{leftmargin=1.2em,itemsep=1pt,topsep=2pt}
\renewcommand{\arraystretch}{1.18}
\newcommand{\bin}{\Delta f}
\pagestyle{empty}
\begin{document}

{\LARGE\bfseries\color{ac} FFT \& Signal Playground}\hfill{\large Formula sheet}\\[-2pt]
{\small\color{gray} Every quantity below is the one computed by \texttt{signals.py}. Symbols: $f_s$ sample rate, $N$ number of samples, $x[n]$ signal, $w[n]$ window, $k$ frequency bin.}

\section{Sampling}
\begin{align*}
T &= \frac{N}{f_s} \qquad
\bin = \frac{f_s}{N} \qquad
f_{\mathrm{Nyq}} = \frac{f_s}{2} \qquad
f_k = k\,\bin \quad (k = 0,\dots,N/2)\\[2pt]
f_{\mathrm{alias}} &= \bigl|\,f - f_s \,\mathrm{round}(f/f_s)\,\bigr| \qquad (\text{a tone with } f > f_{\mathrm{Nyq}} \text{ appears at } f_{\mathrm{alias}})
\end{align*}
Resolving two tones needs a separation of at least about the main-lobe width of the window: $|f_2 - f_1| \gtrsim 2\bin$ for a Hann window. Finer resolution needs a longer record (larger $N$), not a higher $f_s$.

\section{Discrete Fourier transform}
\begin{align*}
X[k] &= \sum_{n=0}^{N-1} x[n]\,e^{-j 2\pi k n/N}
\qquad\qquad
x[n] = \frac{1}{N}\sum_{k=0}^{N-1} X[k]\,e^{\,j 2\pi k n/N}\\
\sum_{n} |x[n]|^2 &= \frac{1}{N}\sum_{k} |X[k]|^2 \quad(\text{Parseval})
\qquad
X[N-k] = X[k]^{*}\ \text{for real } x
\end{align*}
The real-input FFT keeps only $k = 0,\dots,N/2$. An FFT costs $O(N\log N)$ operations against $O(N^2)$ for the direct sum.

\section{Windowing and amplitude scaling}
The spectrum of the windowed signal is the true spectrum convolved with the window's transform. Define the \emph{coherent gain}
\[
\mathrm{CG} = \frac{1}{N}\sum_{n} w[n]
\qquad\Longrightarrow\qquad
A_k = \frac{2\,|X_w[k]|}{N\cdot\mathrm{CG}} \quad (1 \le k < N/2)
\]
so a sine of amplitude $A$ on a bin centre reads exactly $A$ for every window. The DC and Nyquist bins are not doubled. Levels in the plots are
$20\log_{10}A_k$ (dB re amplitude 1). The equivalent noise bandwidth is
\[
\mathrm{ENBW} = N\,\frac{\sum_n w[n]^2}{\bigl(\sum_n w[n]\bigr)^2}\ \text{bins}
\]

\begin{center}
\begin{tabular}{@{}l l c c c c c@{}}
\toprule
Window & $w[n]$ for $0 \le n < N$ & CG & ENBW & Null at & Peak sidelobe & Worst scallop\\
\midrule
Rectangular & $1$ & 1.000 & 1.00 & $\pm 1$ bin & $-13.3$ dB & $-3.90$ dB\\
Hann & $0.5 - 0.5\cos\frac{2\pi n}{N}$ & 0.500 & 1.50 & $\pm 2$ bins & $-31.5$ dB & $-1.42$ dB\\
Hamming & $0.54 - 0.46\cos\frac{2\pi n}{N}$ & 0.540 & 1.36 & $\pm 2$ bins & $-42.7$ dB & $-1.75$ dB\\
Blackman & $0.42 - 0.5\cos\frac{2\pi n}{N} + 0.08\cos\frac{4\pi n}{N}$ & 0.420 & 1.73 & $\pm 3$ bins & $-58.1$ dB & $-1.10$ dB\\
\bottomrule
\end{tabular}
\end{center}
\emph{Scalloping} is the amplitude error when a tone falls halfway between two bins (last column). A narrower main lobe (rectangular) gives the best resolution but the worst leakage. A wider lobe (Blackman) hides weak tones next to strong ones.

\section{Peak frequency by parabolic interpolation}
With the largest bin $k$ and its neighbours $\alpha = A_{k-1}$ and $\beta = A_k$ and $\gamma = A_{k+1}$:
\[
\delta = \frac{1}{2}\,\frac{\alpha - \gamma}{\alpha - 2\beta + \gamma}
\qquad
f_{\mathrm{peak}} = (k + \delta)\,\bin
\qquad
A_{\mathrm{peak}} \approx \beta - \tfrac{1}{4}(\alpha - \gamma)\,\delta
\qquad (-\tfrac12 \le \delta \le \tfrac12)
\]
This locates a tone to a small fraction of a bin (about $0.1$ Hz in the app at $\bin \approx 2$ Hz).

\clearpage
\section{Filters}
Magnitude responses used in the app ($n$ is the filter order $f_c$ the cutoff):
\begin{align*}
|H_{\mathrm{LP}}(f)| &= \frac{1}{\sqrt{1 + (f/f_c)^{2n}}}
\qquad\quad
|H_{\mathrm{HP}}(f)| = \frac{1}{\sqrt{1 + (f_c/f)^{2n}}}\\[2pt]
|H_{\mathrm{BP}}(f)| &= \frac{1}{\sqrt{1 + u^{2n}}}
\qquad
|H_{\mathrm{BS}}(f)| = \frac{1}{\sqrt{1 + u^{-2n}}}
\qquad
u = \frac{f^2 - f_0^2}{f\,B}
\quad f_0 = \sqrt{f_1 f_2}\quad B = f_2 - f_1
\end{align*}
\begin{itemize}
\item $|H| = 1/\sqrt{2}$ ($-3$ dB) at every cutoff. The roll-off is $20n$ dB per decade ($6n$ dB per octave).
\item The filter is applied as $Y[k] = H(f_k)\,X[k]$ on a copy zero-padded to the next power of two $\ge 2N$, then cropped back to $N$ samples. Because $H$ is real the output has \textbf{zero phase shift}, and the padding keeps circular wrap-around out of the result.
\item Convolution in time is multiplication in frequency: $\;y = x * h \;\Longleftrightarrow\; Y = X\,H$.
\end{itemize}

\section{Spectrogram (short-time Fourier transform)}
\[
S[m,k] = \sum_{n=0}^{L-1} x[n + mH]\,w[n]\,e^{-j2\pi kn/L}
\qquad
\mathrm{level}[m,k] = 20\log_{10}\frac{|S[m,k]|}{\sum_n w[n]/2}
\]
with a Hann window of length $L \le 256$ and hop $H = L/4$. Frequency resolution is $f_s/L$ and time resolution is about $L/f_s$, so their product is fixed: a longer window sharpens frequency and blurs time. Values are clipped at $-80$ dB.

\section{Useful theorems and pairs}
\begin{center}
\begin{tabular}{@{}l l@{\qquad}l l@{}}
\toprule
Property & Statement & Pair & Transform\\
\midrule
Linearity & $ax + by \leftrightarrow aX + bY$ & Rectangle (width $T$) & $T\,\mathrm{sinc}(fT)$\\
Time shift & $x[n-n_0] \leftrightarrow X[k]\,e^{-j2\pi kn_0/N}$ & Gaussian & Gaussian\\
Modulation & $x[n]\,e^{j2\pi k_0 n/N} \leftrightarrow X[k-k_0]$ & Impulse train & Impulse train\\
Multiplication & $xy \leftrightarrow X * Y\,/\,N$ & Pure sine & two spectral lines\\
\bottomrule
\end{tabular}
\end{center}
The \emph{Morse code} preset is a keyed carrier $m(t)\sin(2\pi f_0 t)$. By the modulation theorem its spectrum is the spectrum of the on/off envelope $m(t)$ shifted to $\pm f_0$. Faster keying makes $m(t)$ change faster, so the sidebands spread wider around the carrier. The 5 ms raised-cosine edges used in the app make the sidebands fall off much faster than hard on/off keying, which is what removes key clicks.

\section{Worked example}
For $f_s = 8000$ Hz and $N = 4096$:
\[
T = 0.512\ \mathrm{s}
\qquad
\bin = 1.953\ \mathrm{Hz}
\qquad
f_{\mathrm{Nyq}} = 4000\ \mathrm{Hz}
\qquad
\text{a 5 kHz tone aliases to } |5000 - 8000| = 3000\ \mathrm{Hz}
\]
A 4th-order low-pass at $f_c = 900$ Hz passes a 500 Hz tone with
$|H| = \bigl(1 + (500/900)^8\bigr)^{-1/2} \approx 0.995$ and a 1300 Hz tone with $|H| = \bigl(1 + (1300/900)^8\bigr)^{-1/2} \approx 0.224$.

\end{document}
